% =============================================================================
% 5.6 CICLO DE VIDA Y CONSUMO ASÍNCRONO DE SERVICIOS
% =============================================================================
\section{Ciclo de Vida y Consumo Asíncrono de Servicios}
\label{sec:servicios}

Todo el consumo de la API REST del backend (capítulo~\ref{ch:backend},
pág.~\pageref{ch:backend}) se canaliza a través de una única instancia de Axios
preconfigurada. Esta sección documenta su configuración, los interceptores que inyectan el
JWT y gestionan las sesiones expiradas, y el catálogo de servicios de dominio.

\subsection{Configuración del cliente base (\texttt{apiClient.ts})}
\label{subsec:apiclient}

El cliente HTTP se centraliza en \comando{apiClient.ts}, que todos los servicios de
\pathbreak{shared/services} reutilizan. Sus parámetros clave figuran en la
tabla~\ref{tab:apiclient}.

\vspace{0.2cm}
\noindent
\renewcommand{\arraystretch}{1.3}
\fontsize{8.5}{10.2}\selectfont\sffamily
\begin{tabularx}{\textwidth}{|L{3.4cm}|L{3.2cm}|Y|}
\hline
\rowcolor{headerTabla}
\sffamily\textbf{Parámetro} & \sffamily\textbf{Valor} & \sffamily\textbf{Motivo} \\
\hline
\texttt{baseURL} & \pathbreak{/api} o \texttt{VITE\_API\_URL} & Prefijo común; configurable por entorno. \\ \hline
\rowcolor{grisTecnico}
\texttt{timeout} & \texttt{8000} ms & Aborta peticiones colgadas y libera el indicador de carga. \\ \hline
\texttt{headers} & \texttt{Content-Type: application/json} & Contrato JSON por defecto. \\ \hline
\rowcolor{grisTecnico}
\texttt{withCredentials} & \texttt{true} & Envía \textit{cookies} de credenciales (token vía \textit{cookie}). \\ \hline
\end{tabularx}
\label{tab:apiclient}

\begin{NotaTecnica}
En desarrollo, Vite expone un \textit{proxy} que reenvía \pathbreak{/api}, \pathbreak{/auth}
y \pathbreak{/v1} a \pathbreak{http://localhost:8080} (el backend Spring), evitando
problemas de CORS durante el desarrollo local. En producción no hay \textit{proxy}: el
mismo \textit{JAR} sirve la SPA y la API bajo el mismo origen.
\end{NotaTecnica}

La figura~\ref{fig:pipeline-cliente} muestra el recorrido completo de una petición desde
una vista hasta el backend y de vuelta, atravesando los interceptores que se documentan en
los apartados siguientes. Es el espejo, en el cliente, del \textit{pipeline} de servidor
de la figura~\ref{fig:pipeline-http} (pág.~\pageref{fig:pipeline-http}).

\vspace{0.3cm}
\begin{figure}[h!]
\centering
\begin{tikzpicture}[node distance=0.5cm and 0.85cm,
    etapa/.style={rectangle, rounded corners=3pt, draw=c4Sistema!70!black, fill=c4Componente!30,
        font=\sffamily\scriptsize, align=center, minimum width=2.2cm, minimum height=0.8cm, inner sep=2pt},
    intc/.style={etapa, fill=colorAcento!25, draw=colorAcento!80!black},
    rel/.style={-{Stealth[length=2mm]}, semithick, draw=grisISO}]
    \node[etapa, fill=c4Persona!25] (view) {Vista / hook};
    \node[etapa, right=of view] (svc) {Servicio\\\scriptsize de dominio};
    \node[intc, right=of svc] (reqi) {Interceptor\\\scriptsize petición (JWT)};
    \node[etapa, fill=c4Datos!22, right=of reqi] (be) {Backend\\\scriptsize REST};
    \node[intc, below=0.8cm of be] (respi) {Interceptor\\\scriptsize respuesta};
    \node[etapa, fill=c4Persona!25, left=of respi] (store) {Store / UI};

    \draw[rel] (view) -- (svc);
    \draw[rel] (svc) -- (reqi);
    \draw[rel] (reqi) -- node[c4etiqueta, above]{\scriptsize Bearer} (be);
    \draw[rel] (be) -- (respi);
    \draw[rel] (respi) -- node[c4etiqueta, above]{\scriptsize 401/403?} (store);
\end{tikzpicture}
\caption{Pipeline de una petición en el cliente, con sus interceptores.}
\label{fig:pipeline-cliente}
\end{figure}

\subsection{Interceptor de peticiones: inyección del JWT}
\label{subsec:interceptor_req}

El interceptor de petición lee el token (de \texttt{sessionStorage}, con \textit{fallback}
a \texttt{localStorage}) y lo inyecta dinámicamente como cabecera \comando{Authorization:
Bearer <token>}, de modo que ningún servicio gestione el token manualmente.

\begin{lstlisting}[language=JavaScript, caption={Interceptor de petición: inyección del JWT (\texttt{apiClient.ts}).}, label={lst:req-interceptor}]
function readToken(): string | null {
  return sessionStorage.getItem(ACCESS_TOKEN_KEY)
      ?? localStorage.getItem(ACCESS_TOKEN_KEY);
}

apiClient.interceptors.request.use((config) => {
  const token = readToken();
  if (token) config.headers.Authorization = `Bearer ${token}`;
  return config;
});
\end{lstlisting}

\subsection{Interceptor de respuestas: manejo de sesiones expiradas}
\label{subsec:interceptor_resp}

El interceptor de respuesta concentra la \textbf{lógica de detección de sesión expirada}.
Ante un \texttt{401} o \texttt{403}, no asume ciegamente que la sesión caducó:
\textbf{decodifica el \textit{claim} \texttt{exp} del JWT} y lo contrasta con la marca de
expiración almacenada. Solo si el token está efectivamente expirado emite el evento
\comando{auth:unauthorized}, que dispara el cierre de sesión y la redirección. La
figura~\ref{fig:flow-401} formaliza esta decisión.

\vspace{0.3cm}
\begin{figure}[h!]
\centering
\begin{tikzpicture}[node distance=0.75cm and 1.4cm]
    \node[flujoInicio] (resp) {Respuesta HTTP};
    \node[flujoDecision, below=of resp] (d1) {¿401 / 403?};
    \node[flujoProceso, right=2.2cm of d1] (ok) {Propagar al servicio};
    \node[flujoDecision, below=1.0cm of d1] (d2) {¿exp vencido\\o token ausente?};
    \node[flujoProceso, right=2.2cm of d2] (keep) {Mantener sesión\\(error puntual)};
    \node[flujoInicio, fill=bgFail, draw=red!70!black, below=1.0cm of d2] (evt) {Emitir \texttt{auth:unauthorized}};

    \draw[flujoFlecha] (resp) -- (d1);
    \draw[flujoFlecha] (d1) -- node[c4etiqueta, above]{no} (ok);
    \draw[flujoFlecha] (d1) -- node[c4etiqueta, left]{sí} (d2);
    \draw[flujoFlecha] (d2) -- node[c4etiqueta, above]{no} (keep);
    \draw[flujoFlecha] (d2) -- node[c4etiqueta, left]{sí} (evt);
\end{tikzpicture}
\caption{Decisión del interceptor de respuesta ante errores 401/403.}
\label{fig:flow-401}
\end{figure}

\clearpage
\begin{lstlisting}[language=JavaScript, caption={Interceptor de respuesta: distinción entre expiración real y error puntual (\texttt{apiClient.ts}).}, label={lst:resp-interceptor}]
apiClient.interceptors.response.use(
  (response) => response,
  (error: AxiosError) => {
    const status = error.response?.status;
    if (status === 401 || status === 403) {
      const token     = readToken();
      const expiresAt = sessionStorage.getItem(EXPIRES_AT_KEY)
                     ?? localStorage.getItem(EXPIRES_AT_KEY);
      const isExpiredByStore = !token
        || (!!expiresAt && Date.now() > Number(expiresAt) * 1000);
      const isExpiredByJwt   = token ? isJwtExpired(token) : true;
      if (isExpiredByStore || isExpiredByJwt) {
        window.dispatchEvent(new CustomEvent('auth:unauthorized'));
      }
    }
    return Promise.reject(error);
  },
);
\end{lstlisting}

\begin{AvisoNota}
Esta distinción evita cerrar la sesión del usuario por un \texttt{403} legítimo (p.~ej.,
acceder a un recurso ajeno protegido por las políticas RLS del capítulo~\ref{ch:basedatos},
pág.~\pageref{ch:basedatos}) o por un \texttt{401} transitorio: solo se fuerza el
\textit{logout} cuando el JWT está \textbf{realmente} caducado, mejorando la experiencia y
evitando expulsiones espurias. El interceptor también libera el indicador de carga global
ante \textit{timeouts} (\texttt{ECONNABORTED}) y errores de servidor ($\ge$500).
\end{AvisoNota}

\subsection{Catálogo de servicios de dominio}
\label{subsec:servicios_catalogo}

La tabla~\ref{tab:servicios-fe} relaciona los servicios de \pathbreak{shared/services} con
su responsabilidad; todos construidos sobre \texttt{apiClient} y devolviendo \textit{tipos}
de TypeScript derivados del contrato \texttt{ApiResponse<T>} del backend.

\vspace{0.2cm}
\noindent
\renewcommand{\arraystretch}{1.3}
\fontsize{8.5}{10.2}\selectfont\sffamily
\begin{tabularx}{\textwidth}{|L{4.4cm}|Y|}
\hline
\rowcolor{headerTabla}
\sffamily\textbf{Servicio} & \sffamily\textbf{Responsabilidad} \\
\hline
\texttt{authService} & Registro, \textit{login}, sincronización OAuth, refresco de sesión. \\
\hline
\rowcolor{grisTecnico}
\texttt{cvStudioService} & Operaciones del CV: edición, asistencia IA, compilación a PDF. \\
\hline
\texttt{dashboardService} & Métricas y analíticas del reclutador. \\
\hline
\rowcolor{grisTecnico}
\texttt{educationService} · \texttt{experienceService} & Formación y experiencia laboral (con geolocalización). \\
\hline
\texttt{fileService} & Subida de archivos e integración con Drive. \\
\hline
\rowcolor{grisTecnico}
\texttt{notificationsService} · \texttt{preferencesService} & Notificaciones y preferencias del usuario. \\
\hline
\texttt{portfolioService} · \texttt{companyProfileService} & Portafolio público y perfil de empresa. \\
\hline
\rowcolor{grisTecnico}
\texttt{adminService} & Operaciones del panel de administración. \\
\hline
\end{tabularx}
\label{tab:servicios-fe}

\clearpage
\subsection{Resiliencia del cliente y manejo de errores}
\label{subsec:resiliencia}

El cliente aplica tres mecanismos de resiliencia que, combinados, evitan que un error
puntual degrade toda la experiencia. La tabla~\ref{tab:resiliencia} los resume.

\vspace{0.2cm}
\noindent
\renewcommand{\arraystretch}{1.3}
\fontsize{8.5}{10.2}\selectfont\sffamily
\begin{tabularx}{\textwidth}{|L{3.6cm}|Y|}
\hline
\rowcolor{headerTabla}
\sffamily\textbf{Mecanismo} & \sffamily\textbf{Función} \\
\hline
\texttt{RouteErrorBoundary} & Captura errores de renderizado por ruta y muestra una vista de error en lugar de una pantalla en blanco. \\ \hline
\rowcolor{grisTecnico}
\texttt{Suspense} + \textit{skeleton} & Mantiene la interfaz responsiva mientras carga el \textit{chunk} diferido de una página. \\ \hline
\texttt{AbortController} & Cancela peticiones en vuelo al desmontar la vista, evitando condiciones de carrera. \\ \hline
\rowcolor{grisTecnico}
\textit{timeout} 8\,s + \textit{toast} & Aborta peticiones colgadas y notifica al usuario con un \textit{toast} (Sonner). \\ \hline
\end{tabularx}
\label{tab:resiliencia}

\paragraph{Cancelación de peticiones.} El módulo exporta \comando{createAbortController()},
un \textit{factory} de \texttt{AbortController} que las vistas usan para \textbf{cancelar}
peticiones en vuelo al desmontarse (p.~ej., al navegar fuera de una página de búsqueda),
evitando actualizaciones de estado sobre componentes ya desmontados y condiciones de
carrera. El interceptor de respuesta distingue además las cancelaciones
(\texttt{isCancel}) de los errores reales, de modo que abortar una petición no se
contabiliza como fallo.

\paragraph{Separación de responsabilidades en el consumo.} El diseño del consumo de
servicios respeta una separación de responsabilidades nítida que conviene subrayar como
cierre. La \textbf{vista} (o su \textit{custom hook}) decide \textit{cuándo} pedir datos y
\textit{qué} hacer con ellos, pero no conoce el detalle del transporte. El \textbf{servicio
de dominio} traduce esa intención a una llamada concreta y devuelve tipos de TypeScript, sin
saber nada de tokens ni de sesiones. El \textbf{interceptor} inyecta el JWT y gobierna las
sesiones expiradas de forma transversal, sin que ningún servicio lo gestione. Y el
\textbf{store} mantiene el estado resultante, notificando a las vistas suscritas. Esta
estratificación —vista, servicio, interceptor, store— es la contraparte cliente de las
capas Controller~/~Service~/~Repository del backend (capítulo~\ref{ch:backend},
sección~\ref{sec:clases}, pág.~\pageref{sec:clases}): en ambos extremos, cada capa tiene una
única razón de cambio, lo que hace el sistema predecible y mantenible a lo largo de su
ciclo de vida, tal como persigue la norma ISO/IEC/IEEE~15289:2024.
